% 第15章第4节：加权逼近与广义Rayleigh商（广义Rayleigh商的极值、加权低秩逼近）。

\section{加权逼近与广义Rayleigh商}
\label{sec:ls-generalized-rayleigh}

前两节都用欧氏长度：单位方向按$\bm u^\top\bm u=1$规定，重建误差按各元素平方和计算。各坐标的单位或重要性不同时，长度往往按第~\ref{chap:inner-product-spaces}章的加权内积度量。度量一变，有两处跟着改变：规定“单位方向”的长度变了，比较重建误差的标准也变了。

本节先处理前一处：保持二次型不变，改用加权长度规定单位方向，得到广义Rayleigh商及其极值。再处理后一处：在行、列权重可以分离的误差下研究低秩逼近，并看到它的最佳方向正是由广义Rayleigh商选出的。广义Rayleigh商还为后面的图谱松弛提供工具：卷九的谱划分把平衡约束写成按顶点度加权的长度，再用广义Rayleigh商的极值给出离散划分问题的松弛下界，可参考\BookChapterRef{chap:spectral-partition}{卷九“谱划分”章}。

\subsection{广义Rayleigh商的极值}
\label{sec:ls-generalized-eigen-reduction}
\label{sec:ls-generalized-eigen-check}

保持$\bm S=\operatorname{diag}(4,1)$不变，改用$\bm B=\operatorname{diag}(9,1)$规定长度，即要求
\[
 \bm x^\top\bm B\bm x=9x_1^2+x_2^2=1 .
\]
沿第一坐标方向，满足约束的向量是$\pm\bm e_1/3$；沿第二坐标方向，满足约束的向量仍是$\pm\bm e_2$。两者的二次型值分别为
\[
 \Bigl(\frac{\bm e_1}3\Bigr)^{\!\top}\bm S\,\frac{\bm e_1}3=4\cdot\Bigl(\frac13\Bigr)^2=\frac49,\qquad
 \bm e_2^\top\bm S\bm e_2=1 .
\]
在欧氏约束下最优的$\bm e_1$方向，现在不如$\bm e_2$方向。

可见改变长度标准会改变答案。而且与示例~\ref{ex:ls-rayleigh-diagonal-two-by-two}之前的讨论一样，比较两个坐标方向还不够，必须比较全部满足加权约束的方向。与欧氏情形相同，先把约束化为一个不随缩放改变的商。
\begin{definition}[广义Rayleigh商]
\label{def:ls-generalized-rayleigh-quotient}
设$\bm S\in\mathbb R^{d\times d}$对称，$\bm B\in\mathbb R^{d\times d}$对称正定。对$\bm x\neq\bm0$，称
\begin{equation}
 R_{\bm S,\bm B}(\bm x)=\frac{\bm x^\top\bm S\bm x}{\bm x^\top\bm B\bm x}
 \label{eq:ls-generalized-rayleigh-quotient}
\end{equation}
为$\bm S$关于$\bm B$的\term{广义Rayleigh商}（generalized Rayleigh quotient）。
\end{definition}
$\bm B$正定保证分母$\bm x^\top\bm B\bm x$对一切$\bm x\neq\bm0$为正，商有定义。把$\bm x$换成$c\bm x$（$c\neq0$）时分子分母同乘$c^2$，商不变。

对任意$\bm x\neq\bm0$，令$\bm u=\bm x/\sqrt{\bm x^\top\bm B\bm x}$，则
\[
 \bm u^\top\bm B\bm u=1,\qquad R_{\bm S,\bm B}(\bm x)=\bm u^\top\bm S\bm u ;
\]
反过来，约束$\bm x^\top\bm B\bm x=1$上的二次型值就是广义Rayleigh商。所以求广义Rayleigh商的极值，等价于在$\bm x^\top\bm B\bm x=1$上比较$\bm x^\top\bm S\bm x$。

$\bm B=\bm I$时，广义Rayleigh商恰好回到普通Rayleigh商。两者比较的是同一个二次型$\bm x^\top\bm S\bm x$，差别只在于用哪种长度规定“单位输入”：普通Rayleigh商用欧氏长度，广义Rayleigh商用加权长度$\lVert\bm x\rVert_{\bm B}$。

\subsubsection{主要思想}

\paragraph{把加权长度换成欧氏长度}
加权约束$\bm x^\top\bm B\bm x=1$的困难在于它不是欧氏单位球面，第~\ref{sec:ls-rayleigh-extremal}节的结论不能直接使用。思路是换一个坐标，使加权长度在新坐标中成为欧氏长度。

先看开头的例子。约束可以改写成两个平方之和：
\[
 9x_1^2+x_2^2=(3x_1)^2+x_2^2 .
\]
令$z_1=3x_1$、$z_2=x_2$，约束就成为欧氏单位圆$z_1^2+z_2^2=1$。二次型也用新坐标表示，由$x_1=z_1/3$得
\[
 4x_1^2+x_2^2=\frac49z_1^2+z_2^2 .
\]
于是加权约束下的问题，变成了对角矩阵$\operatorname{diag}(4/9,1)$在欧氏单位圆上的普通Rayleigh商问题。

一般的$\bm B$不是对角矩阵，但同样可以写成平方和。设$\bm L$是上三角矩阵，令$\bm z=\bm L\bm x$。二阶时按元素写出
\[
 \begin{bmatrix}z_1\\z_2\end{bmatrix}
 =\begin{bmatrix}l_{1,1}&l_{1,2}\\0&l_{2,2}\end{bmatrix}
 \begin{bmatrix}x_1\\x_2\end{bmatrix}
 =\begin{bmatrix}l_{1,1}x_1+l_{1,2}x_2\\l_{2,2}x_2\end{bmatrix},
\]
展开平方和得
\[
 \begin{aligned}
 z_1^2+z_2^2
 &=l_{1,1}^2x_1^2+2l_{1,1}l_{1,2}x_1x_2+(l_{1,2}^2+l_{2,2}^2)x_2^2\\
 &=\bm x^\top
 \begin{bmatrix}l_{1,1}^2&l_{1,1}l_{1,2}\\l_{1,1}l_{1,2}&l_{1,2}^2+l_{2,2}^2\end{bmatrix}
 \bm x .
 \end{aligned}
\]
中间的矩阵正是$\bm L^\top\bm L$。所以只要$\bm B=\bm L^\top\bm L$，加权长度$\bm x^\top\bm B\bm x$在新坐标中就是欧氏长度$\bm z^\top\bm z$。开头的例子对应$\bm L=\operatorname{diag}(3,1)$。

这样的$\bm L$总是存在。由定理~\ref{thm:ls-gram-factor}，正定的$\bm B$可以写成$\bm B=\bm L^\top\bm L$，其中$\bm L$可逆。于是$\bm x=\bm L^{-1}\bm z$，且$\bm x\neq\bm0$当且仅当$\bm z\neq\bm0$。

\paragraph{化为普通Rayleigh商}
把$\bm x=\bm L^{-1}\bm z$代入分子与分母，得
\[
 \bm x^\top\bm B\bm x=(\bm L\bm x)^\top(\bm L\bm x)=\bm z^\top\bm z,\qquad
 \bm x^\top\bm S\bm x=\bm z^\top\bm L^{-\top}\bm S\bm L^{-1}\bm z ,
\]
其中$\bm L^{-\top}$表示$(\bm L^{-1})^\top$。记
\begin{equation}
 \bm C=\bm L^{-\top}\bm S\bm L^{-1}.
 \label{eq:ls-generalized-rayleigh-reduced-matrix}
\end{equation}
由$\bm S$对称，
\[
 \bm C^\top=\bm L^{-\top}\bm S^\top\bm L^{-1}=\bm C ,
\]
所以$\bm C$仍是实对称矩阵。广义Rayleigh商在新坐标中写成
\begin{equation}
 R_{\bm S,\bm B}(\bm x)=\frac{\bm z^\top\bm C\bm z}{\bm z^\top\bm z}=R_{\bm C}(\bm L\bm x).
 \label{eq:ls-generalized-rayleigh-reduction}
\end{equation}
加权约束问题就化成了$\bm C$的普通Rayleigh商问题，可以直接应用定理~\ref{thm:matrix-analysis-rayleigh-ritz}。

这一化归只用到Gram因子与换元，不需要矩阵平方根、多元微分或Lagrange乘子法。开头的例子中$\bm C=\operatorname{diag}(4/9,1)$，正是上面按元素得到的矩阵。

\paragraph{换回原坐标}
$\bm C$的最优方向在新坐标中，还要换回原坐标。若$\bm C\bm z=\mu\bm z$，令$\bm v=\bm L^{-1}\bm z$，则
\[
 \bm L^{-\top}\bm S\bm v=\mu\bm L\bm v .
\]
两边左乘$\bm L^\top$，得
\[
 \bm S\bm v=\mu\bm L^\top\bm L\bm v=\mu\bm B\bm v .
\]
反过来，若$\bm S\bm v=\mu\bm B\bm v$且$\bm v\neq\bm0$，左乘$\bm L^{-\top}$并令$\bm z=\bm L\bm v$，就得到$\bm C\bm z=\mu\bm z$。所以，$\bm C$的特征对与下面方程的解一一对应。

\begin{definition}[广义特征值]
\label{def:matrix-analysis-generalized-eigenvalue}
对实对称$\bm S$与对称正定$\bm B$，称满足
\begin{equation}
\begin{gathered}
\bm S\bm v=\mu\bm B\bm v\\
\bm v\neq\bm0
\end{gathered}
\label{eq:matrix-analysis-generalized-eigenproblem}
\end{equation}
的$\mu$为广义特征值，$\bm v$为属于$\mu$的广义特征向量，求解它们的问题称为\term{广义特征问题}（generalized eigenvalue problem）。
\end{definition}

广义特征值不依赖Gram因子$\bm L$的选取。由$\bm S-\mu\bm B=\bm L^\top(\bm C-\mu\bm I)\bm L$，
\[
 \det(\bm S-\mu\bm B)=(\det\bm L)^2\det(\bm C-\mu\bm I),
\]
而$\det\bm L\neq0$，所以广义特征值恰是$\det(\bm S-\mu\bm B)=0$的根，也就是$\bm C$的特征值。不同的$\bm L$给出不同的$\bm C$，特征值却相同。

$\bm C$对称，由谱定理，这些根全是实数，按重数计共有$d$个，下面记为$\mu_1\geq\cdots\geq\mu_d$。

\paragraph{极值与最优方向}
把定理~\ref{thm:matrix-analysis-rayleigh-ritz}用于$\bm C$并换回原坐标，就得到加权约束问题的解。

\begin{theorem}[广义Rayleigh商的极值]
\label{thm:ls-generalized-rayleigh-extremal}
设$\bm S$对称、$\bm B$对称正定，广义特征值为$\mu_1\geq\cdots\geq\mu_d$。则对一切$\bm x\neq\bm0$，
\begin{equation}
 \mu_d\leq R_{\bm S,\bm B}(\bm x)\leq\mu_1 .
 \label{eq:ls-generalized-rayleigh-extremal}
\end{equation}
$R_{\bm S,\bm B}(\bm x)=\mu_1$当且仅当$\bm S\bm x=\mu_1\bm B\bm x$，$R_{\bm S,\bm B}(\bm x)=\mu_d$当且仅当$\bm S\bm x=\mu_d\bm B\bm x$。因此，在约束$\bm x^\top\bm B\bm x=1$上，$\bm x^\top\bm S\bm x$的最大值为$\mu_1$，最优解是满足$\bm v^\top\bm B\bm v=1$的属于$\mu_1$的广义特征向量；最小值的情形相同。
\end{theorem}
\begin{proof}
取$\bm B=\bm L^\top\bm L$与$\bm C=\bm L^{-\top}\bm S\bm L^{-1}$如上。$\bm C$的特征值就是$\mu_1\geq\cdots\geq\mu_d$。对$\bm x\neq\bm0$，令$\bm z=\bm L\bm x\neq\bm0$，由$R_{\bm S,\bm B}(\bm x)=R_{\bm C}(\bm z)$及定理~\ref{thm:matrix-analysis-rayleigh-ritz}，得$\mu_d\leq R_{\bm S,\bm B}(\bm x)\leq\mu_1$。

同一定理还说明，$R_{\bm C}(\bm z)=\mu_1$当且仅当$\bm C\bm z=\mu_1\bm z$。由上面的一一对应，这又当且仅当$\bm S\bm x=\mu_1\bm B\bm x$。$\bm C$确有属于$\mu_1$的特征向量，所以上界能够取到。最小值同理。

最后，约束$\bm x^\top\bm B\bm x=1$上的二次型值就是广义Rayleigh商，所以它的最大值为$\mu_1$，且只在满足$\bm S\bm x=\mu_1\bm B\bm x$与$\bm x^\top\bm B\bm x=1$的$\bm x$处取到。
\end{proof}

\paragraph{\texorpdfstring{$\bm B$}{B}正交的广义特征向量}
上面只用到$\bm C$两端的特征向量。对$\bm C$应用完整的谱分解，还能得到全部广义特征向量的结构：它们可以选得关于加权内积$\langle\cdot,\cdot\rangle_{\bm B}$互相正交。

\begin{theorem}[$\bm B$正交的广义特征向量]
\label{thm:ls-generalized-eigen-b-orthogonal}
设$\bm S$对称、$\bm B$对称正定，广义特征值为$\mu_1\geq\cdots\geq\mu_d$。则存在$\bm v_1,\ldots,\bm v_d$，使
\begin{equation}
 \bm S\bm v_i=\mu_i\bm B\bm v_i,\qquad \bm v_i^\top\bm B\bm v_j=\delta_{i,j}\quad(i,j=1,\ldots,d).
 \label{eq:ls-generalized-eigen-b-orthogonal}
\end{equation}
记$\bm V=[\bm v_1,\ldots,\bm v_d]$，$\bm M=\operatorname{diag}(\mu_1,\ldots,\mu_d)$，上式等价于$\bm V^\top\bm B\bm V=\bm I$与$\bm V^\top\bm S\bm V=\bm M$。
\end{theorem}
\begin{proof}
由谱定理，$\bm C=\bm Z\bm M\bm Z^\top$，其中$\bm Z=[\bm z_1,\ldots,\bm z_d]$为正交矩阵，$\bm C\bm z_i=\mu_i\bm z_i$。令$\bm V=\bm L^{-1}\bm Z$，即$\bm v_i=\bm L^{-1}\bm z_i$。由上面的一一对应，$\bm S\bm v_i=\mu_i\bm B\bm v_i$。又
\[
\begin{gathered}
 \bm V^\top\bm B\bm V=\bm Z^\top\bm L^{-\top}\bm L^\top\bm L\bm L^{-1}\bm Z=\bm Z^\top\bm Z=\bm I,\\
 \bm V^\top\bm S\bm V=\bm Z^\top\bm L^{-\top}\bm S\bm L^{-1}\bm Z=\bm Z^\top\bm C\bm Z=\bm M .
\end{gathered}
\]
逐元素比较第一式即得$\bm v_i^\top\bm B\bm v_j=\delta_{i,j}$。
\end{proof}

$\bm V$可逆，所以任意$\bm x$都能写成$\bm x=\sum_ia_i\bm v_i$。由定理，
\[
 \bm x^\top\bm B\bm x=\sum_{i=1}^da_i^2,\qquad \bm x^\top\bm S\bm x=\sum_{i=1}^d\mu_ia_i^2 .
\]
广义Rayleigh商再次成为特征值的加权平均。这与普通情形完全平行，只是坐标改成了$\bm B$正交的广义特征向量。

至此两端极值都已求出。若还需要刻画中间的广义特征值，可以对$\bm C$使用第~\ref{sec:ls-spectral-value-sensitivity}节的Courant--Fischer最小最大原理（定理~\ref{thm:matrix-analysis-courant-fischer}），再按$\bm z=\bm L\bm x$换回原坐标。

\paragraph{用开头的例子核对}
最后回到开头的二维例子，把上面的结论逐一落实，并与直接比较全部可行方向的结果对照。

\begin{example}[加权约束下的最优方向]
\label{ex:ls-generalized-rayleigh-diagonal}
取$\bm S=\operatorname{diag}(4,1)$与$\bm B=\operatorname{diag}(9,1)$，在加权约束$\bm x^\top\bm B\bm x=1$下求$\bm x^\top\bm S\bm x$的最大值、最小值与最优方向。

由$\bm L=\operatorname{diag}(3,1)$得$\bm C=\operatorname{diag}(4/9,1)$。由定理~\ref{thm:matrix-analysis-rayleigh-ritz}，$R_{\bm C}$的最大值为$1$，在$\bm z=\pm\bm e_2$取到；最小值为$4/9$，在$\bm z=\pm\bm e_1$取到。换回原坐标$\bm v=\bm L^{-1}\bm z$，得到
\[
 \text{最大方向}\ \pm\bm e_2,\qquad \text{最小方向}\ \pm\frac{\bm e_1}{3}.
\]

两者都满足加权单位约束：
\[
 \bm e_2^\top\bm B\bm e_2=1,\qquad \Bigl(\frac{\bm e_1}{3}\Bigr)^{\!\top}\bm B\,\frac{\bm e_1}{3}=\frac99=1 .
\]
它们也满足广义特征方程：
\[
 \bm S\bm e_2=\bm e_2=1\cdot\bm B\bm e_2,\qquad
 \bm S\frac{\bm e_1}{3}=\frac43\bm e_1=\frac49\cdot\bm B\frac{\bm e_1}{3},
\]
分别属于$1$与$4/9$；又$\bm e_2^\top\bm B(\bm e_1/3)=0$，满足$\bm B$正交。

不经化归，也可以直接比较全部可行方向。约束给出$x_2^2=1-9x_1^2$，其中$x_1^2\in[0,1/9]$，于是
\[
 \bm x^\top\bm S\bm x=4x_1^2+1-9x_1^2=1-5x_1^2\in\Bigl[\frac49,1\Bigr].
\]
最大值$1$在$x_1=0$即$\pm\bm e_2$处取到，最小值$4/9$在$x_1^2=1/9$即$\pm\bm e_1/3$处取到，与化归的结果一致。

对照示例~\ref{ex:ls-rayleigh-diagonal-two-by-two}：同一个$\bm S$在欧氏约束下的最大方向是$\bm e_1$。加权约束让沿$\bm e_1$的可行向量缩短为$\bm e_1/3$，最大方向就换成了$\bm e_2$。
\end{example}

\subsubsection{算法实现}

上面的例子中$\bm L$是对角矩阵，可以直接求逆。一般情形下，计算也按化归的顺序进行：先求Gram因子，再形成$\bm C$，对$\bm C$作谱分解，最后换回原坐标；换元通过因子分解与三角求解完成，不必显式形成$\bm L^{-1}$。输入是对称的$\bm S$与对称正定的$\bm B$，输出是所需的广义特征值与满足$\bm v^\top\bm B\bm v=1$的广义特征向量。下面用以下新的二维示例逐步演示。

\begin{example}[按三角求解的步骤计算]
\label{ex:matrix-approximation-generalized-cholesky}
取
\[
 \bm S=\begin{bmatrix}2&0\\0&3\end{bmatrix},\qquad
 \bm B=\begin{bmatrix}4&0\\0&1\end{bmatrix},
\]
在约束$\bm x^\top\bm B\bm x=1$下求$\bm x^\top\bm S\bm x$的最大值、最小值与最优方向。
\end{example}

\paragraph{变量}
\begin{itemize}
 \item $\bm L\in\mathbb R^{d\times d}$：$\bm B$的上三角Gram因子，满足$\bm B=\bm L^\top\bm L$。
 \item $\bm W\in\mathbb R^{d\times d}$：中间矩阵$\bm W=\bm L^{-\top}\bm S$。
 \item $\bm C\in\mathbb R^{d\times d}$：化归后的对称矩阵$\bm C=\bm W\bm L^{-1}$。
 \item $\mu_i$与$\bm z_i$：$\bm C$的特征值与单位特征向量，$\mu_i$按降序排列。
 \item $\bm v_i$：换回原坐标后的广义特征向量。
\end{itemize}

\paragraph{输入与预处理}
用第~\ref{sec:ls-cholesky}节的Cholesky分解求出$\bm B=\bm G\bm G^\top$，其中$\bm G$是对角元为正的下三角矩阵，再取
\[
 \bm L=\bm G^\top ,
\]
它是上三角矩阵，满足$\bm B=\bm L^\top\bm L$。若分解中某个开方量不为正，说明$\bm B$不是正定矩阵，正定化归不能继续使用，停止并报告。

例中$\bm B$是对角矩阵，Cholesky分解直接对对角元开方：
\[
 \bm L=\begin{bmatrix}2&0\\0&1\end{bmatrix}.
\]

\paragraph{计算步骤}
\begin{enumerate}
 \item \textbf{形成$\bm W$。}解三角方程组
 \begin{equation}
  \bm L^\top\bm W=\bm S .
  \label{eq:ls-generalized-solve-w}
 \end{equation}
 $\bm L^\top$是下三角矩阵，对$\bm S$的每一列按式~\eqref{eq:ls-forward-substitution}作前代。

 例中$\bm L^\top=\operatorname{diag}(2,1)$，每个方程只含一个未知元：$2w_{1,1}=2$，$w_{2,2}=3$，非对角元的方程右端为零。于是
 \[
  \bm W=\begin{bmatrix}1&0\\0&3\end{bmatrix}.
 \]

 \item \textbf{形成$\bm C$。}由$\bm C=\bm W\bm L^{-1}$得$\bm C\bm L=\bm W$，转置后成为
 \begin{equation}
  \bm L^\top\bm C^\top=\bm W^\top ,
  \label{eq:ls-generalized-solve-c}
 \end{equation}
 仍是以下三角矩阵为系数的方程组，再作一次前代。$\bm C$对称，解出的$\bm C^\top$就是$\bm C$；舍入误差会使算出的两半略有差别，通常取$(\bm C+\bm C^\top)/2$再进入下一步，以保持对称结构。

 例中$2c_{1,1}=1$，$c_{2,2}=3$，于是
 \[
  \bm C=\begin{bmatrix}1/2&0\\0&3\end{bmatrix}.
 \]

 \item \textbf{对$\bm C$作谱分解。}对对称矩阵$\bm C$作正交谱分解，按降序取出特征值与单位特征向量；只需要最大值或最小值时，只取相应的一对。例中
 \[
  \mu_1=3,\quad \bm z_1=\bm e_2;\qquad \mu_2=\frac12,\quad \bm z_2=\bm e_1 .
 \]

 \item \textbf{换回原坐标。}对每个所需的$\bm z_i$解三角方程组
 \begin{equation}
  \bm L\bm v_i=\bm z_i ,
  \label{eq:ls-generalized-recover-v}
 \end{equation}
 $\bm L$是上三角矩阵，按式~\eqref{eq:ls-back-substitution}作回代。由$\bm v_i^\top\bm B\bm v_i=\bm z_i^\top\bm z_i=1$，所得$\bm v_i$自动满足加权单位约束。

 例中$\bm L\bm v_1=\bm e_2$与$\bm L\bm v_2=\bm e_1$分别给出
 \[
  \bm v_1=\begin{bmatrix}0\\1\end{bmatrix},\qquad \bm v_2=\begin{bmatrix}1/2\\0\end{bmatrix}.
 \]

 \item \textbf{核验。}在原坐标中计算广义特征残差$\bm S\bm v_i-\mu_i\bm B\bm v_i$与加权长度$\bm v_i^\top\bm B\bm v_i$，确认前者接近零、后者接近$1$。例中
 \[
  \bm S\bm v_1=\begin{bmatrix}0\\3\end{bmatrix}=3\,\bm B\bm v_1,\qquad
  \bm S\bm v_2=\begin{bmatrix}1\\0\end{bmatrix}=\frac12\begin{bmatrix}2\\0\end{bmatrix}=\frac12\,\bm B\bm v_2 ,
 \]
 两个残差都为零。加权长度为
 \[
  \bm v_1^\top\bm B\bm v_1=1,\qquad \bm v_2^\top\bm B\bm v_2=\Bigl(\frac12\Bigr)^2\cdot4=1 .
 \]
\end{enumerate}

最终得到加权约束下的最大值、最小值与最优方向。例中
\[
 \begin{aligned}
 \max_{\bm x^\top\bm B\bm x=1}\bm x^\top\bm S\bm x&=3, &&\text{在}\ \pm\begin{bmatrix}0\\1\end{bmatrix}\ \text{取到};\\
 \min_{\bm x^\top\bm B\bm x=1}\bm x^\top\bm S\bm x&=\frac12, &&\text{在}\ \pm\begin{bmatrix}1/2\\0\end{bmatrix}\ \text{取到}.
 \end{aligned}
\]

与欧氏约束比较：$\bm S$的普通Rayleigh商最大为$3$、最小为$2$，最优直线同样是两条坐标轴。这里加权没有改变最优直线，却把最小值由$2$变成$1/2$，因为沿$\bm e_1$满足加权约束的向量缩短为$\bm e_1/2$。与示例~\ref{ex:ls-generalized-rayleigh-diagonal}合起来看，长度标准既可能改变最优方向，也可能只改变最优值。

\subsubsection{算法分析}

\paragraph{运算量}
形成$\bm C$只用到Cholesky分解与以三角矩阵为系数的求解，代价都是$O(d^3)$次运算；之后的主要代价在对称特征值计算，所用算法本书不作分析。只需要最大或最小的广义特征对时，也必须先完成同样的换元，因为$\bm S$本身的特征向量一般不是广义特征向量。

\paragraph{为什么不改求普通特征问题}
另一种看似直接的做法是把$\bm S\bm v=\mu\bm B\bm v$左乘$\bm B^{-1}$，改求普通特征问题
\[
 \bm B^{-1}\bm S\bm v=\mu\bm v .
\]
由$\bm B^{-1}=\bm L^{-1}\bm L^{-\top}$，
\[
 \bm B^{-1}\bm S=\bm L^{-1}\bm L^{-\top}\bm S=\bm L^{-1}\bm C\bm L ,
\]
它与$\bm C$相似，特征值相同；但$\bm B^{-1}\bm S$一般不对称，谱定理不能用于它。

看一个具体的例子。取
\[
 \bm S=\begin{bmatrix}2&1\\1&2\end{bmatrix},\qquad
 \bm B=\begin{bmatrix}2&0\\0&1\end{bmatrix},
\]
则
\[
 \bm B^{-1}\bm S=\begin{bmatrix}1&1/2\\1&2\end{bmatrix},
\]
它不对称。它的特征多项式与特征值为
\[
 \mu^2-3\mu+\frac32=0,\qquad \mu_{1,2}=\frac{3\pm\sqrt3}{2}.
\]
由第一行$(1-\mu)x_1+\tfrac12x_2=0$得$x_2=2(\mu-1)x_1$，取$x_1=1$，特征向量为
\[
 \bm v_1=\begin{bmatrix}1\\1+\sqrt3\end{bmatrix},\qquad
 \bm v_2=\begin{bmatrix}1\\1-\sqrt3\end{bmatrix}.
\]
这两个向量就是广义特征向量。它们关于$\bm B$正交，在欧氏内积下却不正交：
\[
 \bm v_1^\top\bm B\bm v_2=2+(1-3)=0,\qquad
 \bm v_1^\top\bm v_2=1+(1-3)=-1\neq0 .
\]
所以不能对$\bm B^{-1}\bm S$套用谱定理与欧氏正交特征基的结论。前面两个二维例子里$\bm B^{-1}\bm S$恰好是对角矩阵，只是特殊情形。

\paragraph{正定性的作用}
本小节的化归依赖$\bm B$正定。若$\bm B$不是正定矩阵，$\bm x^\top\bm B\bm x$不再定义长度：$\bm B$半正定而奇异时，存在$\bm x\neq\bm0$使分母为零；$\bm B$不定时，$\bm x^\top\bm B\bm x$可以为负。这时也不存在可逆的$\bm L$使$\bm B=\bm L^\top\bm L$：若有这样的$\bm L$，则$\bm x\neq\bm0$时$\bm L\bm x\neq\bm0$，
\[
 \bm x^\top\bm B\bm x=\lVert\bm L\bm x\rVert_2^2>0 ,
\]
$\bm B$就是正定的。

于是化归无从进行，约束问题甚至可能没有极值。例如$\bm S=\bm I$、$\bm B=\operatorname{diag}(1,0)$时，约束$x_1^2=1$不限制$x_2$，而
\[
 \bm x^\top\bm S\bm x=1+x_2^2
\]
没有上界。这类更一般的广义特征问题不在此展开。

\subsection{加权低秩逼近}
\label{sec:matrix-approximation-weighted-low-rank}

\subsubsection{主要思想}

\paragraph{加权误差}
现在改变重建误差的度量。设$\bm A\in\mathbb R^{m\times d}$的各行表示不同的量，重要性或单位不同；各列表示不同的样本，可信程度也不同。给定对称正定的行权重$\bm\Omega_{\mathrm r}\in\mathbb R^{m\times m}$与列权重$\bm\Omega_{\mathrm c}\in\mathbb R^{d\times d}$，用
\begin{equation}
 E(\bm X)=\operatorname{tr}\bigl((\bm A-\bm X)^\top\bm\Omega_{\mathrm r}(\bm A-\bm X)\bm\Omega_{\mathrm c}\bigr)
 \label{eq:matrix-approximation-weighted-error}
\end{equation}
度量$\bm X$的误差，并在$\operatorname{rank}\bm X\leq k$下求其最小值，称为\term{加权低秩逼近}（weighted low-rank approximation）。$\bm\Omega_{\mathrm r}=\bm I$、$\bm\Omega_{\mathrm c}=\bm I$时它就是上一节的Frobenius平方误差。

先看两个权重都是对角矩阵的情形，记
\[
 \bm\Omega_{\mathrm r}=\operatorname{diag}(\omega_1,\ldots,\omega_m),\qquad
 \bm\Omega_{\mathrm c}=\operatorname{diag}(\tau_1,\ldots,\tau_d),
\]
并记$\bm M=\bm A-\bm X$。矩阵$\bm M^\top\bm\Omega_{\mathrm r}\bm M$的第$(j,j)$元是
\[
 (\bm M^\top\bm\Omega_{\mathrm r}\bm M)_{j,j}=\sum_{i=1}^m\omega_im_{i,j}^2 ;
\]
右乘$\bm\Omega_{\mathrm c}$后它乘上$\tau_j$，再对$j$求迹，得
\[
 E(\bm X)=\sum_{i=1}^m\sum_{j=1}^d\omega_i\tau_j(a_{i,j}-x_{i,j})^2 .
\]
$m=d=2$时按元素写出是
\[
 \begin{aligned}
 E(\bm X)={}&\omega_1\tau_1(a_{1,1}-x_{1,1})^2+\omega_1\tau_2(a_{1,2}-x_{1,2})^2\\
 &+\omega_2\tau_1(a_{2,1}-x_{2,1})^2+\omega_2\tau_2(a_{2,2}-x_{2,2})^2 .
 \end{aligned}
\]
即第$(i,j)$个元素的平方偏差乘以权重$\omega_i\tau_j$。所以这一类误差包含逐元素加权，但要求权重能分离成行因子与列因子的乘积。

\paragraph{化为普通的Frobenius误差}
与上一小节一样，用Gram因子把加权换成普通长度。对角情形先按元素看：取$\bm K=\operatorname{diag}(\sqrt{\omega_1},\ldots,\sqrt{\omega_m})$与$\bm N=\operatorname{diag}(\sqrt{\tau_1},\ldots,\sqrt{\tau_d})$，则
\[
 (\bm K\bm M\bm N)_{i,j}=\sqrt{\omega_i}\,m_{i,j}\sqrt{\tau_j},\qquad
 (\bm K\bm M\bm N)_{i,j}^2=\omega_i\tau_jm_{i,j}^2 ,
\]
对全部元素求和，正好是上面的$E(\bm X)$。

一般情形同样成立。由定理~\ref{thm:ls-gram-factor}，存在可逆的$\bm K$与$\bm N$，使
\[
 \bm\Omega_{\mathrm r}=\bm K^\top\bm K,\qquad \bm\Omega_{\mathrm c}=\bm N\bm N^\top
\]
（对$\bm\Omega_{\mathrm c}$取$\bm N^\top$为它的Gram因子）。由迹的循环性，
\[
 \begin{aligned}
 \lVert\bm K\bm M\bm N\rVert_{\mathrm F}^2
 &=\operatorname{tr}(\bm N^\top\bm M^\top\bm K^\top\bm K\bm M\bm N)\\
 &=\operatorname{tr}(\bm M^\top\bm\Omega_{\mathrm r}\bm M\bm N\bm N^\top)
 =E(\bm X).
 \end{aligned}
\]
于是
\begin{equation}
 E(\bm X)=\lVert\bm K(\bm A-\bm X)\bm N\rVert_{\mathrm F}^2 ,
 \label{eq:matrix-approximation-weighted-frobenius}
\end{equation}
加权误差是变换后矩阵的普通Frobenius误差。

\paragraph{秩约束不变}
令
\begin{equation}
 \widetilde{\bm A}=\bm K\bm A\bm N,\qquad \bm Z=\bm K\bm X\bm N,
 \label{eq:matrix-approximation-weighted-transform}
\end{equation}
则$E(\bm X)=\lVert\widetilde{\bm A}-\bm Z\rVert_{\mathrm F}^2$。

还要确认秩约束在这个变换下不变。乘积的列空间含于左因子的列空间，所以$\operatorname{rank}\bm Z\leq\operatorname{rank}\bm X$；又$\bm X=\bm K^{-1}\bm Z\bm N^{-1}$，同理$\operatorname{rank}\bm X\leq\operatorname{rank}\bm Z$，故二者相等。

于是$\bm X\mapsto\bm K\bm X\bm N$把秩不超过$k$的矩阵一一地变为秩不超过$k$的矩阵，原问题等价于对$\widetilde{\bm A}$作普通的低秩逼近。由定理~\ref{thm:ls-best-low-rank}与定理~\ref{thm:ls-best-low-rank-uniqueness}立即得到下面的结论。

\begin{theorem}[可分离权重下的最佳低秩逼近]
\label{thm:matrix-approximation-weighted-low-rank}
设$\bm\Omega_{\mathrm r}=\bm K^\top\bm K$、$\bm\Omega_{\mathrm c}=\bm N\bm N^\top$，$\bm K$、$\bm N$可逆，$\widetilde{\bm A}=\bm K\bm A\bm N$的奇异值为$\tilde\sigma_1\geq\cdots\geq\tilde\sigma_r>0$，其中$r=\operatorname{rank}\widetilde{\bm A}=\operatorname{rank}\bm A$。对$0\leq k<r$，
\begin{equation}
 \min_{\operatorname{rank}\bm X\leq k}E(\bm X)=\sum_{i=k+1}^r\tilde\sigma_i^2 ,
 \label{eq:matrix-approximation-weighted-low-rank-min}
\end{equation}
最小值由$\bm X_k=\bm K^{-1}\widetilde{\bm A}_k\bm N^{-1}$达到，其中$\widetilde{\bm A}_k$是$\widetilde{\bm A}$的截断SVD。若$1\leq k<r$且$\tilde\sigma_k>\tilde\sigma_{k+1}$，最优解唯一。
\end{theorem}
\begin{proof}
由上面的等式，$E(\bm X)=\lVert\widetilde{\bm A}-\bm K\bm X\bm N\rVert_{\mathrm F}^2$，而$\bm X\mapsto\bm K\bm X\bm N$是秩不超过$k$的矩阵集合到自身的一一映射。所以$E$在原集合上的最小值等于$\lVert\widetilde{\bm A}-\bm Z\rVert_{\mathrm F}^2$在$\operatorname{rank}\bm Z\leq k$上的最小值，由定理~\ref{thm:ls-best-low-rank}，它等于$\sum_{i>k}\tilde\sigma_i^2$，在$\bm Z=\widetilde{\bm A}_k$即$\bm X=\bm X_k$处取到。

有间隙时，由定理~\ref{thm:ls-best-low-rank-uniqueness}，最优的$\bm Z$唯一，一一映射保证最优的$\bm X$也唯一。
\end{proof}

\paragraph{最佳方向由广义Rayleigh商选出}
最优解的结构可以用上一小节的广义Rayleigh商看清。只设行权重，取$\bm\Omega_{\mathrm c}=\bm I$、$\bm N=\bm I$，并取$k=1$。变换后的秩一候选写成$\bm Z=\bm z\bm c^\top$，其中$\lVert\bm z\rVert_2=1$。令$\bm w=\bm K^{-1}\bm z$，则原坐标中的候选是$\bm X=\bm w\bm c^\top$，单位约束变为
\[
 \bm w^\top\bm\Omega_{\mathrm r}\bm w=\bm z^\top\bm z=1 .
\]

按第~\ref{sec:ls-rank-one-rayleigh}节，变换后问题的保留量是$\bm z^\top\widetilde{\bm A}\widetilde{\bm A}^\top\bm z$。代入$\bm z=\bm K\bm w$与$\widetilde{\bm A}=\bm K\bm A$，得
\[
 \begin{aligned}
 \bm z^\top\widetilde{\bm A}\widetilde{\bm A}^\top\bm z
 &=\bm w^\top\bm K^\top\bm K\bm A\bm A^\top\bm K^\top\bm K\bm w\\
 &=\bm w^\top\bm\Omega_{\mathrm r}\bm A\bm A^\top\bm\Omega_{\mathrm r}\bm w .
 \end{aligned}
\]
所以最佳方向$\bm w$使广义Rayleigh商$R_{\bm\Omega_{\mathrm r}\bm A\bm A^\top\bm\Omega_{\mathrm r},\bm\Omega_{\mathrm r}}$取最大值，由定理~\ref{thm:ls-generalized-rayleigh-extremal}，它是广义特征问题
\begin{equation}
 \bm\Omega_{\mathrm r}\bm A\bm A^\top\bm\Omega_{\mathrm r}\bm w=\mu\,\bm\Omega_{\mathrm r}\bm w
 \label{eq:matrix-approximation-weighted-direction}
\end{equation}
的最大特征向量。按上一小节的化归，这里的
\[
 \bm C=\bm K^{-\top}(\bm\Omega_{\mathrm r}\bm A\bm A^\top\bm\Omega_{\mathrm r})\bm K^{-1}=\widetilde{\bm A}\widetilde{\bm A}^\top ,
\]
最大广义特征值就是$\tilde\sigma_1^2$。

对应的最优系数与最优近似是
\begin{equation}
 \begin{aligned}
 \bm c&=\widetilde{\bm A}^\top\bm z=\bm A^\top\bm\Omega_{\mathrm r}\bm w,\\
 \bm X_1&=\bm w\bm c^\top=\bm w\bm w^\top\bm\Omega_{\mathrm r}\bm A .
 \end{aligned}
 \label{eq:matrix-approximation-weighted-rank-one-optimum}
\end{equation}
记$\bm P=\bm w\bm w^\top\bm\Omega_{\mathrm r}$。由$\bm w^\top\bm\Omega_{\mathrm r}\bm w=1$，
\[
 \bm P^2=\bm w(\bm w^\top\bm\Omega_{\mathrm r}\bm w)\bm w^\top\bm\Omega_{\mathrm r}=\bm P ,
\]
且$\bm\Omega_{\mathrm r}\bm P=\bm\Omega_{\mathrm r}\bm w\bm w^\top\bm\Omega_{\mathrm r}$对称，所以$\bm P$是关于加权内积$\langle\cdot,\cdot\rangle_{\bm\Omega_{\mathrm r}}$到$\operatorname{span}\{\bm w\}$的正交投影。

加权秩一逼近因而与第~\ref{sec:ls-rank-one-rayleigh}节完全平行，只是长度与投影都换成了按$\bm\Omega_{\mathrm r}$度量。

权重改变了用于排序的奇异方向：决定保留哪些方向的是$\widetilde{\bm A}$的奇异值，而不是$\bm A$的。因此直接截断原矩阵$\bm A$，一般不再是加权意义下的最优解，下面的计算给出具体比较。

\subsubsection{算法实现}

计算按化归的顺序进行：求出两个权重的三角因子，形成$\widetilde{\bm A}$并截断它的SVD，再用三角求解还原到原坐标。输入是$\bm A$、对称正定的$\bm\Omega_{\mathrm r}$与$\bm\Omega_{\mathrm c}$以及目标秩$k$，输出是秩不超过$k$的$\bm X_k$。下面用以下二维示例逐步演示。

\begin{example}[行加权改变保留的方向]
\label{ex:matrix-approximation-weighted-low-rank}
取
\[
 \bm A=\begin{bmatrix}3&0\\0&2\end{bmatrix},\qquad
 \bm\Omega_{\mathrm r}=\begin{bmatrix}1&0\\0&4\end{bmatrix},
\]
列权重取$\bm\Omega_{\mathrm c}=\bm I$，目标秩$k=1$。求加权误差最小的秩一近似，并与直接截断$\bm A$比较。
\end{example}

\paragraph{变量}
\begin{itemize}
 \item $\bm K\in\mathbb R^{m\times m}$：$\bm\Omega_{\mathrm r}$的上三角Gram因子，满足$\bm\Omega_{\mathrm r}=\bm K^\top\bm K$。
 \item $\bm N\in\mathbb R^{d\times d}$：$\bm\Omega_{\mathrm c}$的下三角Cholesky因子，满足$\bm\Omega_{\mathrm c}=\bm N\bm N^\top$。
 \item $\widetilde{\bm A}\in\mathbb R^{m\times d}$：变换后的矩阵$\bm K\bm A\bm N$。
 \item $\widetilde{\bm U}_k,\widetilde{\bm\Sigma}_k,\widetilde{\bm V}_k$：$\widetilde{\bm A}$的前$k$个左奇异向量、奇异值与右奇异向量。
 \item $\bm F\in\mathbb R^{m\times k}$、$\bm H\in\mathbb R^{d\times k}$：还原后的左、右因子。
 \item $\bm X_k\in\mathbb R^{m\times d}$：原坐标中的近似。
\end{itemize}

\paragraph{输入与预处理}
用第~\ref{sec:ls-cholesky}节的Cholesky分解求出$\bm\Omega_{\mathrm r}=\bm G\bm G^\top$，取$\bm K=\bm G^\top$；再求出$\bm\Omega_{\mathrm c}=\bm N\bm N^\top$，$\bm N$直接取Cholesky因子。权重是对角矩阵时，两个因子都直接取各对角元的平方根。若某个开方量不为正，说明权重不是正定矩阵，停止并报告。

例中
\[
 \bm K=\begin{bmatrix}1&0\\0&2\end{bmatrix},\qquad \bm N=\bm I .
\]

\paragraph{计算步骤}
\begin{enumerate}
 \item \textbf{形成$\widetilde{\bm A}$。}按式~\eqref{eq:matrix-approximation-weighted-transform}作两次矩阵乘法
 \[
  \widetilde{\bm A}=\bm K\bm A\bm N .
 \]
 例中第二行乘以$2$，得
 \[
  \widetilde{\bm A}=\begin{bmatrix}3&0\\0&4\end{bmatrix}.
 \]

 \item \textbf{截断SVD。}求$\widetilde{\bm A}$的SVD，保留前$k$个奇异值及其奇异向量，得到
 \[
  \widetilde{\bm A}_k=\widetilde{\bm U}_k\widetilde{\bm\Sigma}_k\widetilde{\bm V}_k^\top .
 \]
 例中奇异值按降序为$\tilde\sigma_1=4$与$\tilde\sigma_2=3$，前者的左、右奇异向量都是$\bm e_2$，后者都是$\bm e_1$。保留第一个，
 \[
  \widetilde{\bm A}_1=4\,\bm e_2\bm e_2^\top=\begin{bmatrix}0&0\\0&4\end{bmatrix}.
 \]
 注意排在前面的是$\bm A$中较短的第二个坐标方向。

 \item \textbf{还原到原坐标。}要算的是$\bm X_k=\bm K^{-1}\widetilde{\bm A}_k\bm N^{-1}$。不显式求逆，而是对两侧的因子分别解三角方程组
 \begin{equation}
  \bm K\bm F=\widetilde{\bm U}_k,\qquad \bm N^\top\bm H=\widetilde{\bm V}_k ,
  \label{eq:matrix-approximation-weighted-recover-factors}
 \end{equation}
 $\bm K$与$\bm N^\top$都是上三角矩阵，各作一次回代。于是$\bm F=\bm K^{-1}\widetilde{\bm U}_k$，$\bm H^\top=\widetilde{\bm V}_k^\top\bm N^{-1}$，近似以因子形式给出：
 \[
  \bm X_k=\bm F\widetilde{\bm\Sigma}_k\bm H^\top .
 \]
 中间的$\widetilde{\bm\Sigma}_k$是$k\times k$矩阵，秩不超过$k$由因子形式直接保证。

 若需要完整的矩阵而不是因子，也可以先解$\bm N^\top\bm Y^\top=\widetilde{\bm A}_k^\top$得$\bm Y=\widetilde{\bm A}_k\bm N^{-1}$，再解$\bm K\bm X_k=\bm Y$。

 例中$\bm K\bm F=\bm e_2$给出$\bm F=\bm e_2/2$，$\bm N=\bm I$给出$\bm H=\bm e_2$，于是
 \[
  \bm X_1=\frac{\bm e_2}{2}\cdot4\cdot\bm e_2^\top=\begin{bmatrix}0&0\\0&2\end{bmatrix}.
 \]

 \item \textbf{核验。}在原坐标下核验秩不超过$k$，并计算加权残差，与定理~\ref{thm:matrix-approximation-weighted-low-rank}给出的下界$\sum_{i>k}\tilde\sigma_i^2$比较。

 例中$\bm X_1$的秩为$1$，加权残差为
 \[
  \begin{aligned}
  E(\bm X_1)&=\bigl\lVert\bm K(\bm A-\bm X_1)\bm N\bigr\rVert_{\mathrm F}^2\\
  &=\left\lVert\begin{bmatrix}3&0\\0&0\end{bmatrix}\right\rVert_{\mathrm F}^2=9=\tilde\sigma_2^2 ,
  \end{aligned}
 \]
 与下界相等。
\end{enumerate}

最终得到加权意义下的最佳秩一近似。例中
\[
 \bm X_1=\begin{bmatrix}0&0\\0&2\end{bmatrix},\qquad E(\bm X_1)=9 .
\]
直接截断$\bm A$得到$\operatorname{diag}(3,0)$，它的加权误差是
\[
 \left\lVert\bm K\begin{bmatrix}0&0\\0&2\end{bmatrix}\right\rVert_{\mathrm F}^2
 =\left\lVert\begin{bmatrix}0&0\\0&4\end{bmatrix}\right\rVert_{\mathrm F}^2=16 ,
\]
更大。

再用广义Rayleigh商核对选出的方向。例中
\[
 \bm\Omega_{\mathrm r}\bm A\bm A^\top\bm\Omega_{\mathrm r}=\begin{bmatrix}9&0\\0&64\end{bmatrix},
\]
关于$\bm\Omega_{\mathrm r}=\operatorname{diag}(1,4)$的广义特征值为$9/1=9$与$64/4=16$。最大者$16=\tilde\sigma_1^2$对应
\[
 \bm w=\frac{\bm e_2}{2},\qquad \bm w^\top\bm\Omega_{\mathrm r}\bm w=\frac14\cdot4=1 ,
\]
它正是还原步骤中的$\bm F$。按式~\eqref{eq:matrix-approximation-weighted-rank-one-optimum}，系数与重建为
\[
 \bm c=\bm A^\top\bm\Omega_{\mathrm r}\bm w=\begin{bmatrix}0\\4\end{bmatrix},\qquad
 \bm w\bm c^\top=\begin{bmatrix}0&0\\0&2\end{bmatrix}=\bm X_1 .
\]
第二行的权重放大了它的偏差，所以虽然$\bm A$的第一列更长，最优近似却保留了第二个坐标方向。

\subsubsection{算法分析}

\paragraph{运算量与正定性}
除一次$m\times d$矩阵的SVD外，其余步骤是两次Cholesky分解、两次矩阵乘法与若干三角求解，代价为$O(m^3+d^3+md(m+d))$次运算。权重必须正定；某行或某列权重为零时，相应的$\bm K$或$\bm N$不可逆，还原步骤无法进行。

\paragraph{还原步骤的误差放大}
还原步骤会放大计算误差。若算出的$\widetilde{\bm A}_k$带有误差$\bm\Delta$，还原后的误差为$\bm K^{-1}\bm\Delta\bm N^{-1}$。左乘时逐列应用谱范数的定义，记$\bm\Delta$的第$j$列为$\bm\delta_j$，得
\[
 \begin{aligned}
 \lVert\bm K^{-1}\bm\Delta\rVert_{\mathrm F}^2&=\sum_j\lVert\bm K^{-1}\bm\delta_j\rVert_2^2
 \leq\lVert\bm K^{-1}\rVert_2^2\sum_j\lVert\bm\delta_j\rVert_2^2\\
 &=\lVert\bm K^{-1}\rVert_2^2\lVert\bm\Delta\rVert_{\mathrm F}^2 .
 \end{aligned}
\]
右乘时对各行同样处理，合起来是
\begin{equation}
 \lVert\bm K^{-1}\bm\Delta\bm N^{-1}\rVert_{\mathrm F}\leq\lVert\bm K^{-1}\rVert_2\lVert\bm\Delta\rVert_{\mathrm F}\lVert\bm N^{-1}\rVert_2 .
 \label{eq:matrix-approximation-weighted-recovery-error}
\end{equation}
例中$\bm K^{-1}=\operatorname{diag}(1,1/2)$、$\bm N^{-1}=\bm I$，两个谱范数都是$1$，还原不放大误差。若行权重改为$\operatorname{diag}(1,10^{-4})$，则$\bm K^{-1}=\operatorname{diag}(1,100)$，放大倍数可达$100$。

权重的最大、最小特征值相差悬殊时，原坐标下的结果需要直接计算加权残差，并与理论下界$\tilde\sigma_{k+1}^2+\cdots$比较来核验。

\paragraph{权重不能分离时}
这一化归要求误差能写成左右乘法的形式，任意的逐元素权重未必满足。若对一切$\bm X$都有
\[
 E(\bm X)=\sum_{i,j}w_{i,j}(a_{i,j}-x_{i,j})^2 ,
\]
取$\bm A-\bm X=\bm e_i\bm e_j^\top$，左端为$(\bm\Omega_{\mathrm r})_{i,i}(\bm\Omega_{\mathrm c})_{j,j}$，右端为$w_{i,j}$。于是权重矩阵必是行因子与列因子的乘积，满足
\begin{equation}
 w_{i,j}w_{k,l}=w_{i,l}w_{k,j} .
 \label{eq:matrix-approximation-separable-weights}
\end{equation}
例如$\left[\begin{smallmatrix}1&2\\2&1\end{smallmatrix}\right]$不满足$1\cdot1=2\cdot2$，不能这样化归；缺失元素对应权重为零，也不满足正定性要求。

这些情形下误差不再按任何正交投影分解，没有截断SVD式的闭式解，一般只能用迭代优化求解，并且要分别讨论所得解是否为全局最优。这类问题的通用求解方法由卷五展开，可参考\BookChapterRef{chap:optimization}{卷五第1章“优化问题与解”}。
